\chapter{Metodología de Desarrollo y Planificación}
\label{cap:metodologiaDesarrollo}

Uno de los factores más importantes en el desarrollo de un proyecto es el plan y metodología que se va a llevar a cabo para su desarrollo.
En este capítulo detallamos cómo se ha definido la metodología de nuestra aplicación llamada \textit{ClimbIt} para lograr un 
flujo de trabajo funcional con sus herramientas específicas para la gestión y organización que permite ir avanzar eficientemente en las 
diferentes fases del proyecto.

\section{Metodología de trabajo: Scrumban}

Como metodología de trabajo nos decantamos por utilizar \textbf{Scrumban} debido a que consideramos que aplicar Scrum puro podía resultar 
demasiado estricto especialmente para un equipo de dos personas y en Kanban echamos en falta peridos de trabajo definidos para tener la eficiencia  
el trabajo con presión de fechas de entrega.   

De Scrum nos quedamos con el concepto de \textit{sprints} (iteraciones temporales de desarrollo) de dos semanas aproximadamente y de las técnicas
de especificación de requisitos mediante \textit{Hisotiras de usuario} y \textit{Epics} (estructuras para la definición de casos de uso). 
Por otro lado, de Kanban en lugar de preasignar tareas de forma estricta trabajamos con un sistema \textit{pull} de forma que cada miembro 
del equipo selecciona la tarea a realizar de la columna \textit{To Do} .

% [Imagen del tablero de trabajo]

\section{Proceso de Captura de Requisitos y Definición del Producto}

Antes de iniciar la fase de desarrollo e implementación, determinamos que era fundamental comprender en profundidad los 
problemas a resolver, de manera que no cometiéramos el error de basar la programación en suposiciones no contrastadas. Por lo que 
dedicamos una fase inicial a la definición estratégica del producto, apoyándonos en herramientas de diseño ágil.

\subsection{Lienzo de Propuesta de Valor (Value Proposition Canvas)}

Con el objetivo de empatizar de manera efectiva con el usuario final, construimos un Value Proposition Canvas, 
siendo esta una herramienta que nos permite adoptar la perspectiva del escalador y analizar su contexto dividiéndolo en dos 
bloques principales, el Perfil de cliente y Mapa de valor.

\paragraph{Perfil del cliente}
: Se enfoca exclusivamente en entender al usuario. Te pones en su situación y describes sus 
"trabajos" (lo que intenta resolver), sus "dolores" o frustraciones (los obstáculos y emociones negativas que encuentra) 
y sus "alegrías" o ganancias (los resultados y beneficios que desea obtener).

Se identificó que el usuario escalador busca cumplir "trabajos" funcionales (registrar sesiones, ver pistas), 
sociales (compartir logros) y emocionales (sentirse motivado). Sus principales "frustraciones" (Pains) giran en torno a 
la dificultad del registro manual, la falta de una plataforma centralizada y la desmotivación por no visualizar su progreso. 
Por otro lado, sus "alegrías" (Gains) deseadas son la facilidad para seguir su evolución, recibir feedback y sentirse parte de una 
comunidad competitiva y sana.

\paragraph{Mapa de Valor}
Se enfoca en cómo un producto o servicio crea valor en base a los datos del pefil de cliente. Aquí se describen los 
"productos y servicios" alivian los "dolores" de tus clientes (Pain Relivers) y cómo generan las 
"alegrías" que ellos buscan (Gain creators).

Para responder a este perfil, se determino que los Pain Relievers se enfocan en ofrecer una 
interfaz intuitiva, un registro de progreso automatizado y una plataforma segura y centralizada. Finalmente, sus "creadores de alegrías" 
(Gain Creators) se centran en proporcionar estadísticas detalladas, un sistema de logros motivador y una comparativa social que impulse 
al usuario a mejorar e interactuar con sus compañeros.

% [Imagen de VPC]

\subsection{Visualización del viaje: User Story Map}

A partir de la propuesta de valor, se construyó un User Story Map para visualizar el "viaje" completo del usuario dentro de la aplicación, 
utilizando la plataforma Mural.  Un USR (Mapa de Historias de Usuario) es una técnica visual utilizada en la gestión ágil de proyectos para 
organizar y priorizar las historias de  usuario. Este enfoque facilita la comprensión del flujo de trabajo y las interacciones del usuario 
con el sistema, ayudando a definir y planificar las funcionalidades del producto de manera efectiva.

Este mapa visual resulto muy importante ya que nos ayudó a dividir el proyecto. Trazando líneas horizontales sobre el mapa, 
pudimos decidir qué funcionalidades eran vitales para el primer prototipo, cuáles entrarían en el Producto Mínimo Viable (MVP) y qué ideas 
quedaban relegadas a versiones futuras.

%[Captura del USM]

\section{Herramientas de gestión y seguimiento}

Para la gestión del proyecto tomamos la decisión de tener las herramientas lo más cercanas posibles para reducir el rozamiento de su uso.
 Descartamos usar plataformas externas tipo Jira o Trello y centralizamos toda la gestión directamente en \textbf{GitHub Projects}, 
 manteniendo los requisitos vinculados al código fuente y las herramientas de despliegue. 

El \textit{Product Backlog} se fue rellenando progresivamente con las tareas extraídas del User Story Map. Cada requisito, ya fuera una Historia 
de Usuario (HU) funcional o una tarea técnica (como configurar la base de datos), se tradujo en una \textit{Issue} dentro del repositorio 
con información suficiente para no perder contexto.

\begin{table}[h]
\centering
\begin{tabular}{|l|p{10cm}|}
\hline
\textbf{Atributo} & \textbf{Descripción y uso en el proyecto} \\ \hline
\textbf{Estado} & Columna del tablero Kanban (\textit{Backlog, To Do, In Progress, Review, Done}). \\ \hline
\textbf{Labels} & Categorización visual (\textit{bug, backend, frontend, documentación, epic, HU...}). \\ \hline
\textbf{Estimación} & Tamaño (\textit{Size}) desde XS hasta XL para medir el esfuerzo relativo. \\ \hline
\textbf{Prioridad} & Nivel de urgencia (\textit{P0, P1, P2}), marcando P0 las tareas críticas de bloqueo. \\ \hline
\textbf{Milestone} & Vinculación directa al \textit{sprint} actual para seguimiento temporal. \\ \hline
\textbf{Criterios} & Listado de condiciones exactas para considerar la tarea como finalizada. \\ \hline
\end{tabular}
\caption{Anatomía de configuración de una \textit{Issue} en GitHub Projects.}
\label{tab:config_issues}
\end{table}

Además, el seguimiento de los cuatro \textit{sprints} se apoyó en un \textit{Roadmap} visual y en gráficos de \textit{Burndown}, 
que nos permitía visualizar el progreso y detectar si habíamos sobrestimado nuestra capacidad de desarrollo.

% [capturas de roadmap, y de un burdown chart]

\section{Gestión de la Configuración y Control de Versiones}

Implementamos un flujo de trabajo basado en \textbf{Git Flow}, estableciendo reglas claras sobre cómo se integra el 
trabajo de cada desarrollador.

En el repositorio existen dos ramas principales protegidas. La rama \texttt{main} se reservó exclusivamente para las 
versiones estables y liberadas del producto de manera que todo el código de esta rama se subia a producción, mientras que \texttt{develop} 
actuaba como el entorno de preproducción para los nuevos casos de uso. Cada \textit{Issue} del tablero requería la creación de una rama 
específcia (\textit{feature branch}). Al finalizar el desarrollo, esta rama se integraba en \texttt{develop} mediante una estrategia de 
\textit{squash and merge}. Esto nos permitió agrupar decenas de pequeños \textit{commits} en un solo bloque semántico, manteniendo el 
historial del proyecto limpio y legible.

En paralelo a los ciclos de \textit{sprint}, el despliegue de versiones se gestionó a través de la herramienta de \textit{Releases} 
de GitHub. Adoptamos el versionado semántico clásico (X.Y.Z) y mantenemos la trazabilidad de los cambios mediante un  
archivo \texttt{CHANGELOG.md} que sigue los estándares de \textit{Keep a Changelog}. Además para mantener la buena praxis el repositorio 
cuenta desde el inicio con archivos `.gitignore` configurados para el \textit{stack} tecnológico y una guía de instalación en el `README.md`.

\section{Uso de Inteligencia Artificial como Apoyo al Desarrollo}

(por hacer)